% !TeX program = pdflatex
\documentclass[11pt,a4paper]{amsart}

\newcommand{\notelanguage}{finnish}
% Shared preamble for course notes and exercise sets.

\usepackage[T1]{fontenc}
\usepackage{lmodern}
\providecommand{\notelanguage}{english}
\usepackage[finnish,english]{babel}
\AtBeginDocument{\expandafter\selectlanguage\expandafter{\notelanguage}}

\newcommand{\notesfinnishlanguage}{finnish}
\ifx\notelanguage\notesfinnishlanguage
\newcommand{\notestheoremname}{Lause}
\newcommand{\noteslemmaname}{Lemma}
\newcommand{\notespropositionname}{Propositio}
\newcommand{\notescorollaryname}{Korollaari}
\newcommand{\notesdefinitionname}{Määritelmä}
\newcommand{\notesexamplename}{Esimerkki}
\newcommand{\notesproblemname}{Tehtävä}
\newcommand{\notesexercisename}{Tehtävä}
\newcommand{\notesremarkname}{Huomautus}
\newcommand{\notessolutionname}{Ratkaisu}
\else
\newcommand{\notestheoremname}{Theorem}
\newcommand{\noteslemmaname}{Lemma}
\newcommand{\notespropositionname}{Proposition}
\newcommand{\notescorollaryname}{Corollary}
\newcommand{\notesdefinitionname}{Definition}
\newcommand{\notesexamplename}{Example}
\newcommand{\notesproblemname}{Problem}
\newcommand{\notesexercisename}{Exercise}
\newcommand{\notesremarkname}{Remark}
\newcommand{\notessolutionname}{Solution}
\fi

\usepackage[a4paper,margin=30mm]{geometry}
\usepackage{microtype}
\usepackage{graphicx}
\usepackage{enumitem}

\newcommand{\exerciseheading}[2]{%
  \par\noindent
  {\large\bfseries #1\par}
  \vspace{0.1em}
  \noindent{\large #2\par}
  \vspace{1.5em}
}

\usepackage{amsmath}
\usepackage{amssymb}
\usepackage{amsthm}
\usepackage{mathtools}
\usepackage{aliascnt}

% Row operations beside the affected row, e.g. \rowop{+2R_1}.
\newcommand{\rowop}[1]{\mathrel{\big|}\,#1}

% Linear systems: the argument gives the number of alignment pairs.
\newenvironment{linearsystem}[1]
{\left\{
  \begin{alignedat}{#1}}
    {
  \end{alignedat}\right.}

  \usepackage[hidelinks,pdfusetitle]{hyperref}

  \numberwithin{equation}{section}
  \setlist{itemsep=0.25em,topsep=0.5em}

  \theoremstyle{plain}
  \newtheorem{theorem}{\notestheoremname}[section]
  \newaliascnt{lemma}{theorem}
  \newtheorem{lemma}[lemma]{\noteslemmaname}
  \aliascntresetthe{lemma}
  \newaliascnt{proposition}{theorem}
  \newtheorem{proposition}[proposition]{\notespropositionname}
  \aliascntresetthe{proposition}
  \newaliascnt{corollary}{theorem}
  \newtheorem{corollary}[corollary]{\notescorollaryname}
  \aliascntresetthe{corollary}

  \theoremstyle{definition}
  \newaliascnt{definition}{theorem}
  \newtheorem{definition}[definition]{\notesdefinitionname}
  \aliascntresetthe{definition}
  \newaliascnt{example}{theorem}
  \newtheorem{example}[example]{\notesexamplename}
  \aliascntresetthe{example}
  \newaliascnt{problem}{theorem}
  \newtheorem{problem}[problem]{\notesproblemname}
  \aliascntresetthe{problem}
  \newtheorem{exercise}{\notesexercisename}

  \theoremstyle{remark}
  \newaliascnt{remark}{theorem}
  \newtheorem{remark}[remark]{\notesremarkname}
  \aliascntresetthe{remark}

  \newenvironment{solution}
  {
  \begin{proof}[\notessolutionname]}
    {
  \end{proof}}

  \usepackage[nameinlink,noabbrev]{cleveref}

  \crefname{theorem}{theorem}{theorems}
  \crefname{lemma}{lemma}{lemmas}
  \crefname{proposition}{proposition}{propositions}
  \crefname{corollary}{corollary}{corollaries}
  \crefname{definition}{definition}{definitions}
  \crefname{example}{example}{examples}
  \crefname{problem}{problem}{problems}
  \crefname{exercise}{exercise}{exercises}
  \crefname{remark}{remark}{remarks}

  \ifx\notelanguage\notesfinnishlanguage
  \crefname{theorem}{lause}{lauseet}
  \crefname{lemma}{lemma}{lemmat}
  \crefname{proposition}{propositio}{propositiot}
  \crefname{corollary}{korollaari}{korollaarit}
  \crefname{definition}{määritelmä}{määritelmät}
  \crefname{example}{esimerkki}{esimerkit}
  \crefname{problem}{tehtävä}{tehtävät}
  \crefname{exercise}{tehtävä}{tehtävät}
  \crefname{remark}{huomautus}{huomautukset}
  \fi

  \newcommand{\NN}{\mathbb{N}}
  \newcommand{\NNone}{\mathbb{N}_1}
  \newcommand{\ZZ}{\mathbb{Z}}
  \newcommand{\QQ}{\mathbb{Q}}
  \newcommand{\RR}{\mathbb{R}}
  \newcommand{\CC}{\mathbb{C}}
  \newcommand{\PP}{\mathbb{P}}

  \DeclareMathOperator{\im}{im}
  \DeclareMathOperator{\rank}{rank}
  \DeclareMathOperator{\syt}{syt}
  \DeclareMathOperator{\pyj}{pyj}
  \DeclareMathOperator{\dom}{dom}
  \DeclarePairedDelimiter{\abs}{\lvert}{\rvert}
  \DeclarePairedDelimiter{\norm}{\lVert}{\rVert}
  \newcommand{\set}[1]{\left\lbrace#1\right\rbrace}
  \newcommand{\maxset}[1]{\max\set{#1}}
  \newcommand{\minset}[1]{\min\set{#1}}
  \newcommand{\powerset}[1]{\mathcal{P}\!\left(#1\right)}


\title{Analyysi II\\Tutkijapolku 2}
\author{}
\date{}

\begin{document}

\exerciseheading{Analyysi II}{Tutkijapolku 2}

\begin{exercise}
  Olkoon \(I \subset \RR\) väli ja \(f, g \colon I \to \RR\) konvekseja ja
  \[
    f(x), g(x) \geq 0
  \]
  kaikilla \(x \in I\). Anna esimerkki tilanteesta, jossa \(fg\) ei
  ole konveksi.
\end{exercise}

\begin{solution}
  Olkoon \(I = [0, 1]\) ja määritellään \(f, g \colon I \to \RR\) asettamalla
  \[
    f(x) = x, \qquad g(x) = 1 - x.
  \]
  Molemmat funktiot ovat epänegatiivisia välillä \(I\).
  Olkoot \(x, y \in I\) ja \(\lambda \in [0, 1]\) mielivaltaisia. Tällöin
  \[
    \begin{aligned}
      f\bigl(\lambda x + (1 - \lambda)y\bigr)
      &= \lambda f(x) + (1 - \lambda)f(y), \\
      g\bigl(\lambda x + (1 - \lambda)y\bigr)
      &= \lambda g(x) + (1 - \lambda)g(y),
    \end{aligned}
  \]
  joten \(f\) ja \(g\) ovat konvekseja. Niiden tulolle pätee
  \[
    h(x) = (fg)(x) = x(1 - x).
  \]
  Pisteissä \(x = 0\) ja \(y = 1\) konveksiusehto ei kuitenkaan toteudu
  valinnalla \(\lambda = \tfrac{1}{2}\), sillä
  \[
    h\!\left(\frac{0 + 1}{2}\right)
    = \frac{1}{4}
    > 0
    = \frac{h(0) + h(1)}{2}.
  \]
  Siis \(fg\) ei ole konveksi.
\end{solution}

\begin{exercise}
  Olkoon \(I \subset \RR\) väli ja \(f, g \colon I \to \RR\) konvekseja ja
  \[
    f(x), g(x) \geq 0
  \]
  kaikilla \(x \in I\). Anna jotkin sellaiset ehdot funktioille \(f, g\),
  että \(fg\) on varmasti konveksi. Pohdi, onko antamasi ehto paras
  mahdollinen.
\end{exercise}

\begin{solution}
  Oletetaan tehtävän oletusten lisäksi, että
  \begin{equation}\tag{*}\label{eq:fg-condition}
    \bigl(f(x) - f(y)\bigr)\bigl(g(x) - g(y)\bigr) \geq 0
    \qquad \text{kaikilla } x, y \in I.
  \end{equation}
  Osoitetaan, että tällöin \(fg\) on konveksi.

  Olkoot \(x, y \in I\) ja \(\lambda \in [0, 1]\) mielivaltaisia, ja merkitään
  \(z = (1 - \lambda)x + \lambda y \in I\).
  Funktioiden \(f\) ja \(g\) konveksiuden ja epänegatiivisuuden nojalla
  \[
    \begin{aligned}
      (fg)(z)
      &= f(z)g(z) \\
      &\leq \bigl[(1 - \lambda)f(x) + \lambda f(y)\bigr]
      \bigl[(1 - \lambda)g(x) + \lambda g(y)\bigr] \\
      &= (1 - \lambda)f(x)g(x) + \lambda f(y)g(y) \\
      &\qquad {}- \lambda(1 - \lambda)
      \bigl(f(x) - f(y)\bigr)\bigl(g(x) - g(y)\bigr) \\
      &\leq (1 - \lambda)(fg)(x) + \lambda(fg)(y).
    \end{aligned}
  \]
  Viimeinen epäyhtälö seuraa ehdosta~\eqref{eq:fg-condition}.
  Siis \(fg\) on konveksi.

  \newpage
  \begin{samepage}
    Oletetaan seuraavaksi, että \(f\) ja \(g\) ovat lisäksi jatkuvia
    välillä \(I\) ja kaksi kertaa derivoituvia sen sisäpisteissä.
    Tällöin toinen riittävä ehto tulon konveksiudelle on
    \begin{equation}\tag{**}\label{eq:fg-condition-2}
      2f'(x)g'(x) \geq 0
      \qquad \text{kaikilla } x \in I^\circ.
    \end{equation}
    Ehdon riittävyys seuraa tulon derivoimissäännöstä:
    \[
      (fg)''(x)
      = f''(x)g(x) + 2f'(x)g'(x) + f(x)g''(x).
    \]
    Ensimmäinen ja viimeinen termi ovat epänegatiivisia funktioiden
    \(f\) ja \(g\) konveksiuden ja epänegatiivisuuden nojalla.
    Keskimmäinen termi on epänegatiivinen ehdon~\eqref{eq:fg-condition-2}
    nojalla. Siis \((fg)''(x) \geq 0\) kaikilla \(x \in I^\circ\),
    joten tulon jatkuvuuden nojalla \(fg\) on konveksi koko välillä \(I\).\par
  \end{samepage}

  Näillä lisäoletuksilla ehto~\eqref{eq:fg-condition-2} on aidosti
  yleisempi kuin ehto~\eqref{eq:fg-condition}.
  Jos nimittäin \eqref{eq:fg-condition} pätee, niin kaikilla
  \(x, y \in I\), \(x \ne y\), saadaan
  \[
    \frac{f(y) - f(x)}{y - x}\,
    \frac{g(y) - g(x)}{y - x} \geq 0.
  \]
  Kun \(x \in I^\circ\) ja \(y \to x\), saadaan
  \(f'(x)g'(x) \geq 0\), joten myös
  \eqref{eq:fg-condition-2} pätee.

  \begin{samepage}
    Käänteinen väite ei kuitenkaan päde. Valitaan
    \[
      I = [-1, 1], \qquad f(x) = x^2, \qquad g(x) = x^2(3 + x).
    \]
    Molemmat funktiot ovat epänegatiivisia välillä \(I\) ja konvekseja, sillä
    \[
      f''(x) = 2, \qquad g''(x) = 6(1 + x) \geq 0.
    \]
    Ehto~\eqref{eq:fg-condition-2} toteutuu, koska
    \[
      2f'(x)g'(x) = 12x^2(x + 2) \geq 0
      \qquad (x \in I).
    \]
    Sen sijaan ehto~\eqref{eq:fg-condition} ei toteudu:
    \[
      \begin{aligned}
        \bigl(f(-1) - f(\tfrac34)\bigr)
        \bigl(g(-1) - g(\tfrac34)\bigr)
        &= \left(1 - \frac9{16}\right)
        \left(2 - \frac{135}{64}\right) \\
        &= -\frac{49}{1024} < 0.
      \end{aligned}
    \]
    Siis \eqref{eq:fg-condition} ei ole välttämätön tulon konveksiudelle.\par
  \end{samepage}

  \begin{samepage}
    Myöskään ehto~\eqref{eq:fg-condition-2} ei ole välttämätön.
    Esimerkiksi välillä \(I = [0, 1]\) funktiot
    \[
      f(x) = x^2, \qquad g(x) = 3 - x
    \]
    ovat epänegatiivisia ja konvekseja, mutta
    \(2f'(x)g'(x) = -4x < 0\) kaikilla \(x \in (0, 1)\).
    Silti niiden tulo \(h(x) = 3x^2 - x^3\) on konveksi, sillä
    \[
      h''(x) = 6(1 - x) \geq 0
      \qquad (x \in I).
    \]
    Kumpikin annetuista ehdoista on siis riittävä, mutta ei
    välttämätön. Toisin sanoen, olemme löytäneet kaksi eri ehto funktion
    \(fg\) konveksiudelle, mutta kumpikaan näistä ei ole paras
    mahdollinen. \qedhere\par
  \end{samepage}
\end{solution}

\end{document}
